% This is a template for the assignment reports
% TARI29 Artificial Intelligence, 7.5 credits, Winter 2024
% Original version by:  Vladimir Tarasov, 2024
% Revised by:           Alexandros Tzanetos, 2024 

% Use a modern class instead of an old article class
\documentclass{scrartcl}

\KOMAoptions{
    parskip=half,  % full, off
    fontsize=12pt, % base font size (10pt default)
    % headings=big,% small/normal/big headings (normal is default), 
    % paper=a5,    % paper format (a4 default) 
    pagesize=auto  % Use paper format for PDF too
}

% ---------------------- Report details --------------------- %
\newcommand{\docauthor}{Ben Bekir Ertugrul, Moritz Schlitter}
\newcommand{\courseyear}{2026}
\newcommand{\coursename}{TDIS22 Deep Learning}
\newcommand{\coursecredits}{7.5 credits}
\newcommand{\fullcoursename}{\coursename, \coursecredits}
\newcommand{\termname}{Spring \courseyear}
\newcommand{\groupnumber}{?}
\newcommand{\reportname}{Assignment 1 Report: Cats and Dogs Classification}
\newcommand{\assignmentname}{Group \groupnumber\ Assignment Report}
% ------------------ End of report details ------------------ %



% ---------- Setting up properties of the PDF file ---------- %
\usepackage{hyperref}
\hypersetup{
    colorlinks=true,
    allcolors=blue,
    pdftitle={\reportname},
    pdfauthor={Name1, Name2},
    pdfsubject={\coursename, \termname}
	% pdfpagemode=FullScreen,
	% pdfpagemode= UseOutlines % the default if present
}
% ------------ End of properties of the PDF file ------------ %



% ----------- Setting up the headers and footers ------------ %
\usepackage{scrlayer-scrpage}
\pagestyle{scrheadings}
\KOMAoptions{% 
    headsepline,% line below the header
    plainheadsepline,% also on scrplain
}
\clearscrheadfoot % Erase all current configurations
% inner part of the header [titel_page]{other_pages}
\ihead[\coursename, \courseyear]{\coursename, \courseyear}
% outer part of the header [titel_page]{other_pages}
\ohead{\assignmentname}
% center part of the footer [titel_page]{other_pages}
\cfoot[\textup{Page \pagemark\ of \ztotpages}]{\textup{Page \pagemark\ of \ztotpages}}
% ------------- End of the headers and footers -------------- %



% ---------------- Setting up the title page ---------------- %
\title{\reportname}
\subtitle{\assignmentname}
%\author{\docauthor}
\author{Ben Bekir Ertugrul \and Moritz Schlitter
}
\date{\today}
% ------------------ End of the title page ------------------ %

\usepackage[english]{babel}

% --------------- Packages and custom commands -------------- %
% Specify the input encoding of the source file as UTF8 
\usepackage[utf8]{inputenc}
% The T1 font encoding is an 8-bit encoding and fully supports words containing accented characters
\usepackage[T1]{fontenc}
% Load Latin modern fonts in outline format
\usepackage{lmodern}
% Load better typewriter font
\usepackage{inconsolata}
% For better hyphenation patterns
\usepackage[english]{babel}
% For improved micro-typography
\usepackage{microtype}
% Switch off extra space after a sentence
\frenchspacing
% To be able to iclude pictures
\usepackage{graphicx}
\usepackage{adjustbox}
\usepackage{subcaption}
\usepackage{zref-totpages} % To get the total number of pages
\usepackage{mathtools,amsmath} % for advanced math formulas
 % To highlight source code
\usepackage{minted}
% Default options for all minted environments
\setminted{
           style=default, % for minted envronment
           autogobble, % automatically remove all common leading whitespace
           xleftmargin=10pt, % indentation to add (on the left) before the listing
           numbersep=6pt, % gap between numbers and start of line
           linenos, % enables line numbers
           mathescape % enables the usual math mode inside comments
}
\setmintedinline[c]{style=default} % for minted inlines
% Flexible handling of verbatim text
\usepackage{fancyvrb}
% For well-spaced lines and guidelines in tables
\usepackage{booktabs}
% To typeset algorithms or pseudocode in LaTeX 
\usepackage{algorithm}
\usepackage{algpseudocode}
% For plots and graphs
\usepackage{pgfplots}
\pgfplotsset{width=10cm,compat=1.9}
% for quotes " " with \enquote{text}
\usepackage{csquotes}
% for \begin{enumerate}[resume] to work. resumes counting instead of starting from 1 again.
\usepackage{enumitem}

% To generate placeholder text - can be deleted when you have written real text
\usepackage{lipsum}
% ----------- End of packages and custom commands ----------- %

\begin{document}

\maketitle

% It can be commented out if you do not want the table of contents
\pagebreak
% \tableofcontents
% \pagebreak

% \include{introduction}
% \include{preprocessing}
% \include{models}
% %\include{results}

% \bibliographystyle{ieeetr}
% %\bibliography{references}

% \appendix 
% \include{appendix}
\section{Results}

\begin{figure}[htbp]
\centerline{\includegraphics[width=\columnwidth]{resources/logs.png}}
\caption{Training and validation loss for all experiments.}
\label{fig:results}
\end{figure}

\subsection{Experiment 1}
The first experiment concerns itself with the training of a simple CNN model used for binary classification of cat and dog images. This sets a baseline performance for the following experiments, which make use of transfer learning.

The top left plot in Figure \ref{fig:results} visualizes the training process of this experiment. In epoch 1, the model achieves a training loss of about 0.5855. The loss gradually decreases until the last epoch, where it ends up at 0.1031. The validation loss starts at 0.8795, then reaches 0.2780 in epoch 5, and finally zig-zags with a slight downwards trend before it ends up at 0.1312 in epoch 25.

\subsection{Experiment 2}
The model for experiment 2 explores the application of transfer learning by replacing only the output layer of a pretrained CNN multiclass classification model.

Figure \ref{fig:results} shows the output of experiment 2 in the top right corner. Being trained for 50 epochs in total, the training loss is reduced from 0.4455 to 0.0929. The validation loss starts off even lower with just 0.3355 in the first epoch. Similar to experiment 1, we can see it rapidly moving up and down, though with far smaller peaks. The final loss is 0.1224.


\section{Discussion}
The results of this assignment show the implications of transfer learning.

In experiment 1, the initial loss in both validation and training set were bigger than 0.5. Since the CNN model was initialized with random weights, it took 14 epochs in total before training loss and validation loss could be reduced to less than 0.3. After 25 epochs, the training loss seems to be stable at about 0.1, while the validation loss seems to still fluctuate between 0.24 and 0.13.

Contrarily, both losses in experiment 2 are less than 0.45 in the first epoch
since the weights in all but the last layer have been optimized on a similar dataset before. We can see that the validation loss is also much more stable. The only apparent disadvantage is that the training loss takes more epochs to converge to the same value. The model in experiment 1 achieves a training loss of 0.1031 in epoch 25, whereas the model from experiment 2 only achieves a loss of 0.1299 at that epoch.
\end{document}
